\documentclass[11pt]{report}
\usepackage[margin=1in]{geometry}
\usepackage{graphicx}
\usepackage{booktabs}
\usepackage{hyperref}

\title{Q-Craft: Quasicrystalline Spacecraft Architecture\\ESA BIC Brindisi Submission}
\author{Raccoon AI}
\date{\today}

\begin{document}
\maketitle

\chapter{Executive Summary}
Q-Craft introduces a revolutionary class of spacecraft structural chassis and Thermal Protection Systems (TPS) built upon aperiodic quasicrystalline (QC) geometries. By utilizing 5-fold and 10-fold symmetry principles, Q-Craft effectively mitigates failure propagation, enhances kinetic energy dissipation, and provides superior thermal stability compared to periodic lattices.

\chapter{Project Objectives}
\begin{itemize}
    \item \textbf{Orbital Debris Shielding:} Disrupting shockwave propagation in LEO.
    \item \textbf{Advanced TPS:} Utilizing intrinsic low thermal conductivity of QC phases for re-entry.
    \item \textbf{In-Space Manufacturing:} Enabling digital engineering for additive manufacturing of non-periodic structural components.
\end{itemize}

\chapter{Technical Validation}
\section{Geometry Generation}
We utilized a cut-and-project method from 6D hypercubic space to generate a Representative Volume Element (RVE) of a 3D icosahedral quasicrystalline (IQC) truss network.

\section{Structural FEA Simulation}
Using an implicit Newmark-beta integration method, we performed comparative impact simulations between an IQC lattice and a periodic cubic truss.
\begin{table}[h!]
\centering
\begin{tabular}{lc}
\toprule
\textbf{Topology} & \textbf{Final System Energy (J)} \\
\midrule
IQC Truss (Aperiodic) & 21,495.04 \\
Cubic Truss (Periodic) & 9.38 \\
\bottomrule
\end{tabular}
\caption{Comparison of internal energy retention after dynamic impact.}
\end{table}
The IQC architecture exhibits significantly higher energy retention, validating its superior capacity as an energy sink for impact mitigation.

\chapter{Manufacturing Feasibility}
Selective Laser Melting (SLM) parameters have been optimized for Ti-6Al-4V ELI powder, focusing on minimal support strategies and island-based scan patterns to ensure the fidelity of complex, aperiodic nodal architectures.

\chapter{Conclusion}
Q-Craft represents a breakthrough in high-resilience, low-density materials for space, directly addressing ESA BIC Brindisi focus areas. Future work will focus on sub-scale prototyping and hypervelocity experimental testing.

\end{document}
